\documentclass[10pt,a4paper]{amsart}
\usepackage{amsmath,amssymb,amsthm,graphicx}
% Retained minimal preamble for the Stacks Project chapter varieties.tex.
% Verbatim excerpt of the chapter preamble: Xy-pic setup (native preamble lines 14-18), the native theorem
% declarations with the mathematical macros (native preamble lines 154-217), and the project own no-op
% definitions of the reference/slogan/history display environments (native wrapper.tex lines 1-3) that the
% counted statement uses for its slogan. Copied mechanically; nothing added, removed or redefined.
% For commutative diagrams we use Xy-pic
\usepackage[all]{xy}

% We use 2cell for 2-commutative diagrams.
\xyoption{2cell}
% Theorem environments.
%
\theoremstyle{plain}
\newtheorem{theorem}[subsection]{Theorem}
\newtheorem{proposition}[subsection]{Proposition}
\newtheorem{lemma}[subsection]{Lemma}

\theoremstyle{definition}
\newtheorem{definition}[subsection]{Definition}
\newtheorem{example}[subsection]{Example}
\newtheorem{exercise}[subsection]{Exercise}
\newtheorem{situation}[subsection]{Situation}

\theoremstyle{remark}
\newtheorem{remark}[subsection]{Remark}
\newtheorem{remarks}[subsection]{Remarks}

\numberwithin{equation}{subsection}

% Macros
%
\def\lim{\mathop{\mathrm{lim}}\nolimits}
\def\colim{\mathop{\mathrm{colim}}\nolimits}
\def\Spec{\mathop{\mathrm{Spec}}}
\def\Hom{\mathop{\mathrm{Hom}}\nolimits}
\def\Ext{\mathop{\mathrm{Ext}}\nolimits}
\def\SheafHom{\mathop{\mathcal{H}\!\mathit{om}}\nolimits}
\def\SheafExt{\mathop{\mathcal{E}\!\mathit{xt}}\nolimits}
\def\Sch{\mathit{Sch}}
\def\Mor{\mathop{\mathrm{Mor}}\nolimits}
\def\Ob{\mathop{\mathrm{Ob}}\nolimits}
\def\Sh{\mathop{\mathit{Sh}}\nolimits}
\def\NL{\mathop{N\!L}\nolimits}
\def\CH{\mathop{\mathrm{CH}}\nolimits}
\def\proetale{{pro\text{-}\acute{e}tale}}
\def\etale{{\acute{e}tale}}
\def\QCoh{\mathit{QCoh}}
\def\Ker{\mathop{\mathrm{Ker}}}
\def\Im{\mathop{\mathrm{Im}}}
\def\Coker{\mathop{\mathrm{Coker}}}
\def\Coim{\mathop{\mathrm{Coim}}}

% Boxtimes
%
\DeclareMathSymbol{\boxtimes}{\mathbin}{AMSa}{"02}

%
% Macros for moduli stacks/spaces
%
\def\QCohstack{\mathcal{QC}\!\mathit{oh}}
\def\Cohstack{\mathcal{C}\!\mathit{oh}}
\def\Spacesstack{\mathcal{S}\!\mathit{paces}}
\def\Quotfunctor{\mathrm{Quot}}
\def\Hilbfunctor{\mathrm{Hilb}}
\def\Curvesstack{\mathcal{C}\!\mathit{urves}}
\def\Polarizedstack{\mathcal{P}\!\mathit{olarized}}
\def\Complexesstack{\mathcal{C}\!\mathit{omplexes}}
% \Pic is the operator that assigns to X its picard group, usage \Pic(X)
% \Picardstack_{X/B} denotes the Picard stack of X over B
% \Picardfunctor_{X/B} denotes the Picard functor of X over B
\def\Pic{\mathop{\mathrm{Pic}}\nolimits}
\def\Picardstack{\mathcal{P}\!\mathit{ic}}
\def\Picardfunctor{\mathrm{Pic}}
\def\Deformationcategory{\mathcal{D}\!\mathit{ef}}
\newenvironment{reference}{}{}
\newenvironment{slogan}{}{}
\newenvironment{history}{}{}
\usepackage{cleveref}
\crefname{originaldefinition}{original definition}{original definitions}
\Crefname{originaldefinition}{Original definition}{Original definitions}
\makeatletter
\newlabel{algebra-definition-perfect}{{030Y}{source}{}{sourcepointer-1}{}}
\newlabel{algebra-definition-perfect@cref}{{[originaldefinition][0][]030Y}{[1][0][]source}{}{}{}}
\makeatother
\crefname{originallemma}{original lemma}{original lemmas}
\Crefname{originallemma}{Original lemma}{Original lemmas}
\makeatletter
\newlabel{algebra-lemma-minimal-prime-reduced-ring}{{00EU}{source}{}{sourcepointer-2}{}}
\newlabel{algebra-lemma-minimal-prime-reduced-ring@cref}{{[originallemma][0][]00EU}{[1][0][]source}{}{}{}}
\makeatother
\crefname{originallemma}{original lemma}{original lemmas}
\Crefname{originallemma}{Original lemma}{Original lemmas}
\makeatletter
\newlabel{algebra-lemma-separable-smooth}{{00TV}{source}{}{sourcepointer-3}{}}
\newlabel{algebra-lemma-separable-smooth@cref}{{[originallemma][0][]00TV}{[1][0][]source}{}{}{}}
\makeatother
\crefname{originallemma}{original lemma}{original lemmas}
\Crefname{originallemma}{Original lemma}{Original lemmas}
\makeatletter
\newlabel{lemma-geometrically-reduced-dense-smooth-open}{{056V}{source}{}{sourcepointer-4}{}}
\newlabel{lemma-geometrically-reduced-dense-smooth-open@cref}{{[originallemma][0][]056V}{[1][0][]source}{}{}{}}
\makeatother
\crefname{originallemma}{original lemma}{original lemmas}
\Crefname{originallemma}{Original lemma}{Original lemmas}
\makeatletter
\newlabel{lemma-perfect-reduced}{{020I}{source}{}{sourcepointer-5}{}}
\newlabel{lemma-perfect-reduced@cref}{{[originallemma][0][]020I}{[1][0][]source}{}{}{}}
\makeatother
\crefname{originaldefinition}{original definition}{original definitions}
\Crefname{originaldefinition}{Original definition}{Original definitions}
\makeatletter
\newlabel{morphisms-definition-smooth}{{01V5}{source}{}{sourcepointer-6}{}}
\newlabel{morphisms-definition-smooth@cref}{{[originaldefinition][0][]01V5}{[1][0][]source}{}{}{}}
\makeatother
\title{The smooth locus of a reduced locally algebraic scheme over a perfect field is a dense open and equals the regular locus}
\author{The Stacks Project Authors}
\date{stacks-project:varieties.tex:local-snapshot-20261011; GNU FDL 1.2 or later}
\begin{document}
\maketitle
\noindent\textit{Editorial note.} All delivered mathematical text is the author's own, copied from the pinned Stacks Project chapter file varieties.tex; no statement or proof step was rewritten, reordered or removed.\par\medskip
\noindent\textit{Editorial note.} The author gives two independent arguments for density and labels them (1) and (2) in the source; both are reproduced as written.\par\medskip
\noindent\textit{Editorial note.} The counted spans carry 6 navigation stubs and no bibliography citation.\par\medskip
\noindent\textit{Editorial note.} The preamble is the chapter's retained minimal preamble (sources/stacks-varieties/retained-preamble.tex lines 1-77), a verbatim excerpt of the native preamble.\par\medskip
\noindent\textit{Editorial note.} The label lemma-dense-smooth-open-variety-over-perfect-field is defined only in varieties.tex, so no cross-chapter restatement had to be weighed for this card.\par\medskip
\noindent\textit{Allowed standard prerequisites (editorial).}\begin{itemize}
\item Perfect field and separability of residue field extensions. Algebra, Definition algebra-definition-perfect and Lemma algebra-lemma-separable-smooth, invoked at proof lines 4741-4744. Supplies the equality between the smooth locus and the regular locus.
\item The smooth locus of a morphism of schemes is open. Morphisms, Definition morphisms-definition-smooth, invoked at proof lines 4744-4746. Supplies openness.
\item The local ring of a reduced scheme at a generic point of a component is a field. Algebra, Lemma algebra-lemma-minimal-prime-reduced-ring, invoked at proof lines 4748-4750. Supplies the first density argument.
\item Reducedness over a perfect field implies geometric reducedness, and the dense smooth open criterion. Same chapter, Lemmas lemma-perfect-reduced and lemma-geometrically-reduced-dense-smooth-open, invoked at proof lines 4751-4753. Supplies the second density argument.
\end{itemize}\medskip
\ifdefined\hypertarget\hypertarget{sourcepointer-1}{}\fi
\noindent\textit{Source reference.} Original definition 030Y: sources/stacks-chow/tags.txt, source lines 3965--3965. Pointer outside the counted span, invoked once for context; the load-bearing argument of the counted statement is the author's proof reproduced in full in the proof span.\par\medskip
\ifdefined\hypertarget\hypertarget{sourcepointer-2}{}\fi
\noindent\textit{Source reference.} Original lemma 00EU: sources/stacks-chow/tags.txt, source lines 548--548. Pointer outside the counted span, invoked once for context; the load-bearing argument of the counted statement is the author's proof reproduced in full in the proof span.\par\medskip
\ifdefined\hypertarget\hypertarget{sourcepointer-3}{}\fi
\noindent\textit{Source reference.} Original lemma 00TV: sources/stacks-chow/tags.txt, source lines 1089--1089. Pointer outside the counted span, invoked once for context; the load-bearing argument of the counted statement is the author's proof reproduced in full in the proof span.\par\medskip
\ifdefined\hypertarget\hypertarget{sourcepointer-4}{}\fi
\noindent\textit{Source reference.} Original lemma 056V: sources/stacks-varieties/pinned-original-source.tex, source lines 4694--4694. Pointer outside the counted span, invoked once for context; the load-bearing argument of the counted statement is the author's proof reproduced in full in the proof span.\par\medskip
\ifdefined\hypertarget\hypertarget{sourcepointer-5}{}\fi
\noindent\textit{Source reference.} Original lemma 020I: sources/stacks-varieties/pinned-original-source.tex, source lines 401--401. Pointer outside the counted span, invoked once for context; the load-bearing argument of the counted statement is the author's proof reproduced in full in the proof span.\par\medskip
\ifdefined\hypertarget\hypertarget{sourcepointer-6}{}\fi
\noindent\textit{Source reference.} Original definition 01V5: sources/stacks-chow/tags.txt, source lines 2434--2434. Pointer outside the counted span, invoked once for context; the load-bearing argument of the counted statement is the author's proof reproduced in full in the proof span.\par\medskip
\begin{lemma}
\label{lemma-dense-smooth-open-variety-over-perfect-field}
Let $k$ be a perfect field. Let $X$ be a locally algebraic
reduced $k$-scheme, for example a variety over $k$. Then we have
$$
\{x \in X \mid X \to \Spec(k)\text{ is smooth at }x\} =
\{x \in X \mid \mathcal{O}_{X, x}\text{ is regular}\}
$$
and this is a dense open subscheme of $X$.
\end{lemma}

\begin{proof}
The equality of the two sets follows immediately from
Algebra, Lemma \ref{algebra-lemma-separable-smooth} and the definitions
(see Algebra, Definition \ref{algebra-definition-perfect} for the definition
of a perfect field). The set is open because the set of points where
a morphism of schemes is smooth is open, see
Morphisms, Definition \ref{morphisms-definition-smooth}.
Finally, we give two arguments to see that it is dense:
(1) The generic points of $X$ are in the set as the local rings at
generic points are fields (Algebra, Lemma
\ref{algebra-lemma-minimal-prime-reduced-ring}) hence regular.
(2) We use that $X$ is geometrically reduced by
Lemma \ref{lemma-perfect-reduced} and hence
Lemma \ref{lemma-geometrically-reduced-dense-smooth-open} applies.
\end{proof}

\end{document}
